\section{Traveling wave analysis}
\label{sec:traveling}
%%%%%%%%%%%%%%%%%%%%%%%%%%%%%%%%%%%%%%%%%%%%%%%%%%%%%%%%%%%%%%%%%%%%%%%%%%%%%

Motivated by laboratory experiments~\cite{khoshnevis2014combustion, gargar2015recovery, gargar2014recovery, zheng2017numerical, anderson2022analysis, blackwell1959factors},
it is common to search for mathematical solutions to combustion models in the form of traveling waves~\cite{ozbag2018traveling, da1999traveling, ghazaryan2013gasless, gordon2006travelling, ghazaryan2011stability}.
Traveling wave solutions can be characterized as solutions invariant with respect to translation in space \cite{volpert1994traveling}.
They are solutions of the underlying partial differential equation that preserve their shape while propagating with a certain velocity \cite{ghazaryan2022introduction}.

\begin{definition}
	\label{def:traveling}
	The solution $\left(\hat\theta,\hat S_{\text{y}},\hat S_{\text{o}}\right)\left(\xi\right)=\left(\theta,S_{\text{y}},S_{\text{o}}\right)\left(x,t\right)$,
	depending on the traveling variable $\xi=x-ct$ with constant $c$, of \eqref{eq:14}-\eqref{eq:16} is called a traveling wave solution if it satisfies
	\begin{equation}
		\label{eq:limits}
		\lim_{\xi\to-\infty} (\hat\theta,\hat S_{\text{y}},\hat S_{\text{o}}) = (\theta^L,S_{\text{y}}^{\text{L}},S_{\text{o}}^L)
		\quad \text{and} \quad
		\lim_{\xi\to\infty} (\hat\theta,\hat S_{\text{y}},\hat S_{\text{o}}) = (\theta^R,S_{\text{y}}^R,S_{\text{o}}^{\text{R}}),
	\end{equation}
	where the limits coincide with the Riemann problem conditions defined in \eqref{eq:Riemann_boundary}.
\end{definition}

Mathematically, to prove the existence of the traveling wave solution in the sense of Def. \ref{def:traveling} is equivalent to investigating the existence of the corresponding ODE solution with $\alpha$- and $\omega$-limits defined as in \eqref{eq:limits}; see \cite{guckenheimer2013nonlinear,sotomayor1979liccoes} for details.
In what follows, we drop the hats in the notation for simplicity.
When convenient, we refer to traveling waves as in Def. \ref{def:traveling} as combustion fronts as it is commonly done in combustion theory
\cite{ghazaryan2022introduction}.

\subsection{Traveling wave description}

Changing the problem variables $\left(x,t\right)$ to the traveling
ones $\left(\xi=x-ct,t\right)$ in \eqref{eq:14}-\eqref{eq:16},
the original system is rewritten as follows:
\begin{equation}\label{ec4b}
	-c\diff{\theta}{\xi}+
	a_{1}\diff{\left(S_{\text{o}}\theta\right)}{\xi}=
	b_{1}S_{\text{y}}S_{\text{o}}\Phi\left(\theta\right)-
	\textcolor{red}{\beta\theta},
\end{equation}

\begin{equation}\label{ec5b}
	-c\diff{S_{\text{y}}}{\xi}+
	a_{2}
	\diff{S_{\text{y}}}{\xi}=
	-b_{2}S_{\text{y}}S_{\text{o}}\Phi\left(\theta\right),
\end{equation}

\begin{equation}\label{ec6b}
	-c
	\diff{S_{\text{o}}}{\xi}+
	a_{3}
	\diff{S_{\text{o}}}{\xi}=
	-b_{3}S_{\text{y}}S_{\text{o}}\Phi\left(\theta\right).
\end{equation}

Therefore we obtain the following equations:
\begin{equation}
	\label{PHI}
	\Phi\left(\theta\right)=
	t^{\ast}k_{p}
	\exp
	\left(
	-\frac{E_{\text{r}}}{R\left(T^{\ast}\theta+T_{\text{res}}\right)}
	\right).
\end{equation}

Considering
\begin{equation}
	\label{ec36c}
	c=
	\frac{a_{2}b_{3}S_{\text{y}}^{\text{L}} + a_{3} b_{2} S_{\text{o}}^{\text{R}}}{b_{2} S_{\text{o}}^{\text{R}} + b_{3} S_{\text{y}}^{\text{L}}}.
\end{equation}
and
\begin{equation}
	\label{SyL}
	S_{\text{y}}^{\text{L}}=
	-\frac{b_{2}}{b_{3}}
	\frac{\left(c-a_{3}\right)}{\left(c-a_{2}\right)}S_{\text{o}}^{\text{R}}
\end{equation}
obtain
\begin{equation}
	\label{Sy}
	S_{\text{y}}=
	S_{\text{y}}^{\text{L}}+
	\frac{b_{2}}{b_{3}}
	\frac{\left(c - a_{3}\right)}{\left(c - a_{2}\right)}S_{\text{o}}=
	\left(S_{\text{o}}^{\text{R}}-S_{\text{o}}\right)
	\frac{S_{\text{y}}^{\text{L}}}{S_{\text{o}}^{\text{R}}}
\end{equation}
and the ODE system
\begin{equation}
	\label{ec38c}
	\diff{S_{\text{o}}}{\xi}=
	\frac{b_{3}S_{\text{y}}S_{\text{o}}}{c-a_{3}}
	\Phi\left(\theta\right).
\end{equation}

\begin{equation}
	\label{ec39c}
	\diff{\theta}{\xi}=
	\frac{\beta\theta\left(c-a_{3}\right)+\left(a_{1}b_{3}\theta-b_{1}\left(c-a_{3}\right)\right)S_{\text{y}}S_{\text{o}}\Phi\left(\theta\right)}{\left(c-a_{1}S_{\text{o}}\right)\left(c -a_{3}\right)}.
\end{equation}
Considering
\begin{align*}
	F_{1}\left(S_{\text{o}},\theta\right) & =
	\frac{b_{3}S_{\text{y}}S_{\text{o}}}{c-a_{3}}
	\Phi\left(\theta\right).
	% b_{3}
	% \frac{S^{L}_{y}}{S^{R}_{\text{o}}}
	% \frac{\Phi\left(\theta\right)}{c-a_{3}}\left(S^{R}_{\text{o}}-S_{\text{o}}\right)S_{\text{o}}.
	\\
	F_{2}\left(S_{\text{o}},\theta\right) & =
	\frac{\beta\theta\left(c-a_{3}\right)+\left(a_{1}b_{3}\theta-b_{1}\left(c-a_{3}\right)\right)S_{\text{y}}S_{\text{o}}\Phi\left(\theta\right)}{\left(c-a_{1}S_{\text{o}}\right)\left(c -a_{3}\right)}.
	% b_{1}\frac{1-q\theta}{c-a_{1}S_{\text{o}}}
	% S^{L}_{y}\Phi\left(\theta\right)
	% \left[S^{2}_{\text{o}}-S^{R}_{\text{o}}S_{\text{o}}+f\left(\theta\right)\right].
\end{align*}
We can rewrite the ODE system as
\begin{equation}\label{eq:system}
	\left\{
	\begin{aligned}
		\diff{S_{\text{o}}}{\xi} & =
		F_{1}\left(S_{\text{o}},\theta\right). \\
		\diff{\theta}{\xi}       & =
		F_{2}\left(S_{\text{o}},\theta\right).
	\end{aligned}
	\right.
\end{equation}


% TODO: Write where F_1 and F_2 come from.
% TODO: Agregar lemas  If c > a3 and c > a1SR o , then λ1 FC < 0 and λ2 FC > 0. (saddle point)

\textcolor{blue}{
	% Substituting \eqref{ec5b} into \eqref{ec6b} and integrating in $\xi$ results in
	% \begin{equation}\label{ec9b}
	% (a_{2} - c) b_{3} S_{\text{y}} - (a_{3} - c)b_{2} S_{\text{o}} = K_1,
	% \end{equation}
	% where $K_1 \in \mathbb{R}$ is constant.
	% Applying the limits $\xi \to \pm \infty$ in \eqref{ec9c} and \eqref{ec9b}, yields
	% Applying the limits $\xi \to \pm \infty$ in \eqref{ec9b}, yields
	% \begin{equation}\label{ec10c}
	% b_{2} \theta^L ( c - a_{1} S_{\text{o}}^L) + b_{1} S_{\text{y}}^{\text{L}}(c - a_{2}) = b_{2} \theta^R(c - a_{1} S_{\text{o}}^{\text{R}}) + b_{1} S_{\text{y}}^R(c - a_{2})
	% \end{equation}
	% and
	% \begin{equation}\label{ec23c}
	% (a_{2} - c) b_{3} S_{\text{y}}^{\text{L}} - ( a_{3} - c )b_{2} S_{\text{o}}^L = (a_{2} - c) b_{3} S_{\text{y}}^R - (a_{3} - c)b_{2} S_{\text{o}}^{\text{R}}.
	% \end{equation}
	% Substituting boundary condition \eqref{eq:Riemann_boundary} into \eqref{ec10c} and \eqref{ec23c}, we obtain:
	% Substituting boundary condition \eqref{eq:Riemann_boundary} into \eqref{ec23c}, we obtain:
	% \begin{equation}
	% \label{ec35c}
	%        \theta^L  =  \theta^R \left(\displaystyle\frac{c - a_{1} S_{\text{o}}^{\text{R}}}{c}\right) - \displaystyle\frac{(a_{3} - a_{2})b_{1} S_{\text{o}}^{\text{R}} S_{\text{y}}^{\text{L}}}{ a_{2} b_{3} S_{\text{y}}^{\text{L}} +  a_{3} b_{2} S_{\text{o}}^{\text{R}}},
	% \end{equation}
	% and
	% From \eqref{ec9b} and $\xi \to -\infty$ results:
	% From \eqref{ec6b}, we obtain:
	% \begin{equation}
	% 	%\label{ec14}
	% 	\label{Sop}
	% 	S_{\text{o}}^{\prime}  =  -\displaystyle\frac{b_{3} S_{\text{y}} S_{\text{o}}}{a_{3} - c}\Phi.
	% \end{equation}
}
% Analogously, substituting \eqref{ec5b} into \eqref{ec4b} and integrating yields
% \beq\label{ec8c}
% \fa{\partial}{\partial \xi}\left[b_{2} c \theta - b_{2} a_{1} S_{\text{o}} \theta + b_{1} c S_{\text{y}} - b_{1} a_{2} S_{\text{y}} \right] = 0.
% \eeq
% \textcolor{red}{
% 	\begin{equation}\label{ec8c_moda}
% 		\displaystyle\frac{\partial}{\partial \xi}\left[-b_{2} c \theta + b_{2} a_{1} S_{\text{o}} \theta - b_{1} c S_{\text{y}} + b_{1} a_{2} S_{\text{y}} \right] = \textcolor{red}{-\beta \theta b_{2}.}
% 	\end{equation}
% 	so:}
% \textcolor{blue}{
% 	\begin{equation}\label{ec8c_modb}
% 		\displaystyle\frac{\partial}{\partial \xi}\left[b_{2}(a_{1} S_{\text{o}} - c) \theta - b_{1}(c - a_{2}) S_{\text{y}} \right] = \textcolor{red}{-\beta \theta b_{2}.}
% 	\end{equation}
% }

% Substituting \eqref{ec6b} into \eqref{ec4b} and integrating yields
% \textcolor{red}{
% \begin{equation}\label{ec9c_moda}
% \displaystyle\frac{\partial}{\partial \xi}\left[-b_{3} c \theta + b_{3} a_{1} S_{\text{o}} \theta - b_{1} c S_{\text{o}} + b_{1} a_{3} S_{\text{o}} \right] = \textcolor{red}{-\beta \theta b_{3}.}
% \end{equation}
% so:
% }
% \textcolor{blue}{
% \begin{equation}\label{ec9c_modb}
% \displaystyle\frac{\partial}{\partial \xi}\left[b_{3}(a_{1} S_{\text{o}} - c) \theta - b_{1}(c - a_{3}) S_{\text{o}} \right] = \textcolor{red}{-\beta \theta b_{3}.}
% \end{equation}
% }
% \textcolor{red}{
% From \eqref{ec8c_modb} and \eqref{ec9c_modb} results:
% \begin{equation}
%     \label{ec10}
%     \displaystyle\frac{\partial}{\partial \xi}\left[(a_{1} S_{\text{o}} - c) \theta - \displaystyle\frac{b_{1}}{b_{2}}(c - a_{2}) S_{\text{y}} \right] =  \displaystyle\frac{\partial}{\partial \xi}\left[(a_{1} S_{\text{o}} - c) \theta - \displaystyle\frac{b_{1}}{b_{3}}(c - a_{3}) S_{\text{o}} \right]
% \end{equation}
% so:
% \begin{equation}
%     \label{ec11}
%      \displaystyle\frac{\partial}{\partial \xi}\left[- \displaystyle\frac{b_{1}}{b_{2}}(c - a_{2}) S_{\text{y}} + \displaystyle\frac{b_{1}}{b_{3}}(c - a_{3}) S_{\text{o}} \right] = 0
% \end{equation}
% Applying the limits $\xi \to \pm \infty$ in \eqref{ec11}, yields:
% \begin{equation}
%     \label{ec12}
%     - \displaystyle\frac{b_{1}}{b_{2}}(c - a_{2}) S^L_y + \displaystyle\frac{b_{1}}{b_{3}}(c - a_{3}) S_{\text{o}}^L = - \displaystyle\frac{b_{1}}{b_{2}}(c - a_{2}) S_{\text{y}}^R + \displaystyle\frac{b_{1}}{b_{3}}(c - a_{3}) S_{\text{o}}^{\text{R}}
% \end{equation}
% }

% \textcolor{red}{From \eqref{ec11} results
% \begin{equation}
%     \label{ec13}
%     - \displaystyle\frac{b_{1}}{b_{2}}(c - a_{2}) S_{\text{y}} + \displaystyle\frac{b_{1}}{b_{3}}(c - a_{3}) S_{\text{o}} = - \displaystyle\frac{b_{1}}{b_{2}}(c - a_{2}) S_{\text{y}}^R + \displaystyle\frac{b_{1}}{b_{3}}(c - a_{3}) S_{\text{o}}^{\text{R}}
% \end{equation}
% }
% \begin{equation}\label{ec9c}
% b_{2}  \theta(c - a_{1} S_{\text{o}}) + b_{1}  S_{\text{y}}(c - a_{2}) = K,
% \end{equation}
% where $K \in \mathbb{R}$ is constant. \\

\textcolor{blue}{
	% From \eqref{ec8c_modb} results:
	% $$
	% 	\begin{array}{rcl}
	% 		b_{2} a_{1} (S_{\text{o}}^{\prime} \theta + S_{\text{o}} \theta^{\prime}) - b_{2} c \theta^{\prime} - b_{1}(c - a_{2})S_{\text{y}}^{\prime} & = & - \beta \theta b_{2}                                                   \\
	% 		(b_{2} a_{1} S_{\text{o}} - b_{2} c) \theta^{\prime}                                                                  & = & (-\beta b_{2} - b_{2} a_{1} S_{\text{o}}^{\prime})\theta + b_{1}(c - a_{2})S_{\text{y}}^{\prime}
	% 	\end{array}
	% $$
	% \begin{equation}
	% 	\label{thetap}
	% 	\theta^{\prime} =\displaystyle\frac{ (b_{1} c - b_{1}a_{2})S_{\text{y}}^{\prime} - b_{2} a_{1} S_{\text{o}}^{\prime}\theta - \beta b_{2} \theta }{-b_{2} c + b_{2} a_{1} S_{\text{o}}}
	% \end{equation}
	% From \eqref{Sy} we can obtain $S_{\text{y}}^{\prime}$ and substituting into \eqref{thetap}, results:
	% \begin{equation}
	% 	\label{thetap1}
	% 	\theta^{\prime} =\displaystyle\frac{ b_{1}(c - a_{2}) \textcolor{red}{\displaystyle\frac{b_{2}}{b_{3}}\displaystyle\frac{(c - a_{3})}{(c - a_{2})} S_{\text{o}}^{\prime} } - b_{2} a_{1} S_{\text{o}}^{\prime}\theta - \beta b_{2} \theta }{-b_{2} c + b_{2} a_{1} S_{\text{o}}}
	% \end{equation}
	% If $c \neq a_{2}$ then:
	% \begin{equation}
	% 	\label{thetap2}
	% 	\theta^{\prime} =\displaystyle\frac{\textcolor{red}{\displaystyle\frac{b_{1}}{b_{3}}(c - a_{3}) S_{\text{o}}^{\prime} } - a_{1} \theta S_{\text{o}}^{\prime} - \beta \theta }{a_{1} S_{\text{o}} - c}
	% \end{equation}
	% simplifying the above equation:
	% \begin{equation}
	% 	\label{thetap3}
	% 	\theta^{\prime} =\displaystyle\frac{\textcolor{red}{\left(\displaystyle\frac{b_{1}}{b_{3}}(c - a_{3}) - a_{1} \theta\right) S_{\text{o}}^{\prime}} - \beta \theta }{a_{1} S_{\text{o}} - c}
	% \end{equation}
	% \begin{equation}
	% 	\label{thetap4}
	% 	\theta^{\prime} =\displaystyle\frac{\left(b_{1} (c - a_{3}) - a_{1} b_{3}\theta\right) S_{\text{o}}^{\prime} - b_{3} \beta \theta  }{b_{3}(a_{1} S_{\text{o}} - c)}
	% \end{equation}
	% \begin{equation}
	% 	\label{thetap5}
	% 	\theta^{\prime} =\displaystyle\frac{b_{3} \beta \theta + \left(a_{1} b_{3}\theta -  b_{1} (c - a_{3})\right) S_{\text{o}}^{\prime}  }{b_{3}(c - a_{1} S_{\text{o}})}
	% \end{equation}
	% Substituting \eqref{Sop} into \eqref{thetap5}, results:
	% \begin{equation}
	% 	\label{thetap6}
	% 	\theta^{\prime} =\displaystyle\frac{b_{3} \beta \theta + \left(a_{1} b_{3}\theta -  b_{1} (c - a_{3})\right) \textcolor{red}{\displaystyle\frac{b_{3} S_{\text{y}} S_{\text{o}}}{c - a_{3}} \Phi}  }{b_{3}(c - a_{1} S_{\text{o}})}
	% \end{equation}
	% \begin{equation}
	% 	\label{thetap7}
	% 	\theta^{\prime} =\displaystyle\frac{\beta \theta (c - a_{3}) + \left(a_{1} b_{3}\theta -  b_{1} (c - a_{3})\right) \textcolor{red}{S_{\text{y}} S_{\text{o}} \Phi}  }{(c - a_{1} S_{\text{o}})(c -a_{3})}
	% \end{equation}
}
\textcolor{blue}{
	% \begin{equation}
	% 	\label{ec37c}
	% 	S_{\text{y}} = S_{\text{y}}^{\text{L}} + \displaystyle\frac{b_{2}}{b_{3}}\displaystyle\frac{(c - a_{3})}{(c - a_{2})}S_{\text{o}} = (S_{\text{o}}^{\text{R}} - S_{\text{o}})\displaystyle\frac{S_{\text{y}}^{\text{L}}}{S_{\text{o}}^{\text{R}}},
	% \end{equation}
}
\subsection{Equilibrium points}

To find the traveling wave solution of the system~\eqref{Ec1a}-\eqref{Ec3a}
is equivalent to finding the solution of the system composed of two ODEs~\eqref{ec38c}-\eqref{ec39c}
and one algebraic equation~\eqref{Sy}.
Notice that an equilibrium solution of the ODE corresponds to a solution with a derivative equal to
zero~\cite{guckenheimer2013nonlinear,sotomayor1979liccoes}.
In~\eqref{eq:system} equals to zero.
The equation \eqref{ec39c} implies
\begin{equation*}
	F_{1}\left(S_{\text{o}},\theta\right)=0\implies
	S_{\text{y}}=0\quad\vee\quad
	S_{\text{o}}=0.
\end{equation*}
Consider the values $S_{\text{y}}$ and $S_{\text{o}}$, we obtain
\begin{equation*}
	F_{2}\left(S_{\text{o}},\theta\right)=0\implies
	\theta=0.
\end{equation*}
% In any case, from \eqref{ec38c} we obtain
% \begin{equation*}
% 	\diff{\theta}{\xi}=
% 	\frac{\beta\theta}{c-a_{1}S_{\text{o}}}=0,
% \end{equation*}
% therefore
% \begin{equation*}
% 	\theta=
% 	0.
% \end{equation*}

In this paper we are interested in studying the combustion from and
thus we assume that the reaction does not occur at the boundaries,
i.e., the reaction rate $\Phi$ in~\eqref{PHI} vanish when: there
is no oxygen $(Y=0)$, which we call ``Oxygen Controlled'' case
(OC); there is no fuel $(S_{\text{o}}=0)$, which we call
``Fuel Controlled'' case (FC).
In order to approximate applications, we consider that the left
injection end of the domain is FC, and the production right end of
the domain is OC.
The OC and FC equilibria are given by~\eqref{eq:Riemann_boundary}, then:
\begin{align*}
	\text{OC}:
	\big(S^{L}_{\text{o}},\theta^{L}\big) & =
	\left(0,0\right).                         \\
	\text{FC}:
	\big(S^{R}_{\text{o}},\theta^{R}\big) & =
	\big(S^{R}_{\text{o}},0\big).
\end{align*}

From~\eqref{eq:39}, we obtain:
\setcounter{equation}{47}
\begin{equation}
	g\left(0\right)=
	-b_{1}\left(c-a_{3}\right)\Phi\left(0\right).
\end{equation}
From~\eqref{eq:34}, we obtain:
\begin{equation}
	\Phi\left(0\right)=
	t^{\ast}
	k_{p}
	\exp\left(-\frac{E_{\text{r}}}{RT_{\text{res}}}\right).
\end{equation}
At OC and FC equilibria, the matrix $JF$ has the forms:
\begin{equation}\label{eq:50}
	JF_{\text{OC}}=
	\begin{bmatrix}
		\frac{b_{3}S^{L}_{y}}{c-a_{3}}\Phi\left(0\right) & 0               \\
		-\frac{b_{1}S^{L}_{y}}{c}\Phi\left(0\right)      & \frac{\beta}{c}
	\end{bmatrix}.
\end{equation}
\begin{equation}\label{eq:51}
	JF_{\text{FC}}=
	\begin{bmatrix}
		-\frac{b_{3}S^{L}_{y}}{c-a_{3}}\Phi\left(0\right)                & 0                                     \\[.5\baselineskip]
		\frac{b_{1}S^{L}_{y}}{c-a_{1}S^{R}_{\text{o}}}\Phi\left(0\right) & \frac{\beta}{c-a_{1}S^{R}_{\text{o}}}
	\end{bmatrix}.
\end{equation}
because from~\eqref{eq:35} we obtain:
\begin{equation*}
	S_{\text{y}}
	\big(S^{R}_{\text{o}}\big)=0
	\quad\text{ and }\quad
	\diff{S_{\text{y}}}{S_{\text{o}}}=
	-\frac{S^{L}_{y}}{S^{R}_{\text{o}}}.
\end{equation*}

Defining $e=S^{L}_{y}\Phi\left(0\right)$, then the corresponding
eigenvalues and eigenvectors of~\eqref{eq:50} are:
\begin{equation*}
	\begin{vmatrix}
		JF_{\text{OC}}-\lambda I
	\end{vmatrix}=
	\begin{vmatrix}
		\frac{b_{3}e}{c-a_{3}}-\lambda & 0                       \\
		-\frac{b_{1}e}{c}              & \frac{\beta}{c}-\lambda
	\end{vmatrix}=
	\left(\frac{b_{3}e}{c-a_{3}}-\lambda\right)
	\left(\frac{\beta}{c}-\lambda\right)=
	0.
\end{equation*}
then:
\begin{equation*}
	\lambda^{1}_{\text{OC}}=
	\frac{b_{3}e}{c-a_{3}},\qquad
	\lambda^{2}_{\text{OC}}=
	\frac{\beta}{c}.
\end{equation*}

For $\lambda^{1}_{\text{OC}}$:
\begin{equation*}
	\begin{bmatrix}
		0                 & 0                                      \\
		-\frac{b_{1}e}{c} & \frac{\beta}{c}-\frac{b_{3}e}{c-a_{3}}
	\end{bmatrix}
	\begin{bmatrix}
		v_{1} \\
		v_{2}
	\end{bmatrix}=
	\begin{bmatrix}
		0 \\
		0
	\end{bmatrix}\implies
	v_{2}=1\implies
	v=\begin{bmatrix}
		\frac{
			\left(c-a_{3}\right)\beta-b_{3}ec
		}{b_{1}e\left(c-a_{3}\right)} \\
		1
	\end{bmatrix}.
\end{equation*}

For $\lambda^{2}_{\text{OC}}$:
\begin{equation*}
	\begin{bmatrix}
		\frac{b_{3}e}{c-a_{3}}-\frac{\beta}{c} & 0 \\
		-\frac{b_{1}e}{c}                      & 0
	\end{bmatrix}
	\begin{bmatrix}
		v_{1} \\
		v_{2}
	\end{bmatrix}=
	\begin{bmatrix}
		0 \\
		0
	\end{bmatrix}\implies
	v=\begin{bmatrix}
		0 \\
		1
	\end{bmatrix}.
\end{equation*}

The corresponding eigenvalues and eigenvectors of~\eqref{eq:51}
are:
\begin{equation*}
	\begin{vmatrix}
		JF_{\text{FC}}-\lambda I
	\end{vmatrix}=
	\begin{vmatrix}
		-\frac{b_{3}e}{c-a_{3}}-\lambda        & 0                                             \\[.5\baselineskip]
		\frac{b_{1}e}{c-a_{1}S^{R}_{\text{o}}} & \frac{\beta}{c-a_{1}S^{R}_{\text{o}}}-\lambda
	\end{vmatrix}=
	\left(-\frac{b_{3}e}{c-a_{3}}-\lambda\right)
	\left(\frac{\beta}{c-a_{1}S^{R}_{\text{o}}}-\lambda\right)=
	0.
\end{equation*}
then:
\begin{equation*}
	\lambda^{1}_{\text{FC}}=
	-\frac{b_{3}e}{c-a_{3}},\qquad
	\lambda^{2}_{\text{FC}}=
	\frac{\beta}{c-a_{1}S^{R}_{\text{o}}}.
\end{equation*}

For $\lambda^{1}_{\text{FC}}$:
\begin{equation*}
	\begin{bmatrix}
		0                                      & 0                                                            \\[.5\baselineskip]
		\frac{b_{1}e}{c-a_{1}S^{R}_{\text{o}}} & \frac{\beta}{c-a_{1}S^{R}_{\text{o}}}+\frac{b_{3}e}{c-a_{3}}
	\end{bmatrix}
	\begin{bmatrix}
		v_{1} \\
		v_{2}
	\end{bmatrix}=
	\begin{bmatrix}
		0 \\
		0
	\end{bmatrix}\implies
	v_{1}=0\implies
	v=\begin{bmatrix}
		\frac{-\left[\beta\left(c-a_{3}\right)+b_{1}e\left(c-a_{1}S^{R}_{\text{o}}\right)\right]}{b_{1}e\left(c-a_{3}\right)} \\
		1
	\end{bmatrix}.
\end{equation*}

For $\lambda^{2}_{\text{FC}}$:
\begin{equation*}
	\begin{bmatrix}
		-\frac{b_{3}e}{c-a_{3}}-\frac{\beta}{c-a_{1}S^{R}_{\text{o}}} & 0 \\[.5\baselineskip]
		\frac{b_{1}e}{c-a_{1}S^{R}_{\text{o}}}                        & 0
	\end{bmatrix}
	\begin{bmatrix}
		v_{1} \\
		v_{2}
	\end{bmatrix}=
	\begin{bmatrix}
		0 \\
		0
	\end{bmatrix}\implies
	v=\begin{bmatrix}
		0 \\
		1
	\end{bmatrix}.
\end{equation*}

Analyzing the eigenvalues yields the following lemmas

\begin{lemma}
	For $JF_{\text{OC}}\left(0,0\right)$:
	\begin{itemize}
		\item

		      If $c>a_{3}$, then $\lambda^{1}_{\text{OC}}>0$ and
		      $\lambda^{2}_{\text{OC}}>0$.
		      Furthermore, $\left(0,0\right)$ is a repulsor point.
		\item

		      If $c<a_{3}$, then $\lambda^{1}_{\text{OC}}<0$ and
		      $\lambda^{2}_{\text{OC}}>0$.
		      Furthermore, $\left(0,0\right)$ is a saddle point.
	\end{itemize}
\end{lemma}

\begin{lemma}
	For $JF_{\text{FC}}\left(S^{R}_{\text{o}},0\right)$:
	\begin{itemize}
		\item

		      If $c>a_{3}$ and $c>a_{1}S^{R}_{\text{o}}$, then
		      $\lambda^{1}_{\text{FC}}<0$ and
		      $\lambda^{2}_{\text{FC}}>0$.
		      Furthermore, $\left(S^{R}_{\text{o}},0\right)$ is a saddle point.

		\item

		      If $c>a_{3}$ and $c<a_{1}S^{R}_{\text{o}}$, then
		      $\lambda^{1}_{\text{FC}}<0$ and $\lambda^{2}_{\text{FC}}<0$.
		      Furthermore, $\left(S^{R}_{\text{o}},0\right)$ is an
		      attractor point.

		\item

		      If $c<a_{3}$ and $c>a_{1}S^{R}_{\text{o}}$, then
		      $\lambda^{1}_{\text{FC}}>0$ and $\lambda^{2}_{\text{FC}}>0$.
		      Furthermore, $\left(S^{R}_{\text{o}},0\right)$ is a
		      repulsor point.

		\item

		      If $c<a_{3}$ and $c<a_{1}S^{R}_{\text{o}}$, then
		      $\lambda^{1}_{\text{FC}}>0$ and $\lambda^{2}_{\text{FC}}<0$.
		      Furthermore, $\left(S^{R}_{\text{o}},0\right)$ is a
		      saddle point.
	\end{itemize}
\end{lemma}
